\begin{ouverture}[statistiques]{10}{Moyenne, médiane et étendue}{Partie 3 · Données et hasard}
{Moyenne pondérée · Médiane · Étendue · Tableur}
\objectif{Calculer une moyenne pondérée à partir de données, d’un tableau ou d’un diagramme.}
\objectif{Déterminer et interpréter une médiane et une étendue.}
\objectif{Comparer des séries et utiliser un tableur.}
\prerequis{Effectifs, fréquences, moyenne simple et diagrammes (Cinquième).}
\end{ouverture}
\pagestyle{otMain}\essentiel\fontsize{11}{15}\selectfont
\phantomsection\label{cours-10-1}
\essSection{1}{Effectifs et moyenne pondérée}
\begin{definition}L’effectif d’une valeur est le nombre de fois où elle apparaît ; l’effectif total est la somme des effectifs. La fréquence d’une valeur est son effectif divisé par l’effectif total ; les fréquences de toutes les valeurs ont pour somme 1.\end{definition}
\begin{propriete}La moyenne d’une série est la somme de toutes ses valeurs divisée par l’effectif total. Quand les valeurs ont des effectifs, on additionne les produits valeur $\times$ effectif : on calcule une moyenne pondérée par les effectifs.\end{propriete}
\begin{exemple}Valeurs 2, 5 et 8 d’effectifs 4, 3 et 3 : $\frac{2\times4+5\times3+8\times3}{4+3+3}=\frac{47}{10}=4{,}7$. Une moyenne n’est pas forcément une valeur de la série, ni un nombre entier.\end{exemple}
\begin{savoirfaire}[Organiser avant de calculer]Pour une longue liste de données brutes, construire d’abord un tableau d’effectifs. Garder les unités dans les calculs : 126 livres lus par 30 élèves donnent $\frac{126\text{ livres}}{30\text{ élèves}}=4{,}2$ livres par élève en moyenne.\end{savoirfaire}
\begin{attention}On ne calcule pas la moyenne des seules valeurs 2, 5 et 8 sans tenir compte de leurs effectifs. Pour comparer deux groupes de tailles différentes, on compare des fréquences ou des moyennes, pas des totaux.\end{attention}
\clearpage
\phantomsection\label{cours-10-2}
\essSection{2}{Médiane et étendue}
\begin{definition}Une médiane d’une série est une valeur qui partage la série ordonnée en deux parties de même effectif. Pour un effectif impair, c’est la valeur centrale. Pour un effectif pair, toute valeur comprise entre les deux valeurs centrales convient ; on choisit souvent leur moyenne.\end{definition}
\begin{definition}L’étendue d’une série est la différence entre sa plus grande et sa plus petite valeur. Elle renseigne sur la dispersion des valeurs.\end{definition}
\begin{exemple}Pour 2, 4, 4, 7, 8 : la médiane est la 3e valeur, 4, et l’étendue vaut $8-2=6$. Pour 2, 4, 7, 9 : toute valeur comprise entre 4 et 7 est une médiane, par exemple 5,5 ; deux valeurs sont en dessous et deux au-dessus.\end{exemple}
\begin{propriete}[Valeur extrême]Ajouter ou modifier une valeur très grande ou très petite change fortement la moyenne et l’étendue, mais peu la médiane : celle-ci dépend du rang des valeurs, pas de leur taille.\end{propriete}
\begin{savoirfaire}[Avec un tableur]Pour des valeurs saisies dans les cellules A2 à A31, on obtient la moyenne avec \texttt{=MOYENNE(A2:A31)}, une médiane avec \texttt{=MEDIANE(A2:A31)}, l’étendue avec \texttt{=MAX(A2:A31)-MIN(A2:A31)} et l’effectif total avec \texttt{=NB(A2:A31)}.\end{savoirfaire}

\clearpage\rubrique{Découverte}
\begin{blocLivre}{ex-10-1}\begin{exercice}[1]{Un effectif total}
Les notes 8, 10 et 12 ont respectivement pour effectifs 3, 5 et 2. Donner l’effectif total et la fréquence de la note 10.
\end{exercice}\end{blocLivre}
\begin{blocLivre}{ex-10-2}\begin{exercice}[1]{Une moyenne simple}
Calculer la moyenne de 6, 9, 10 et 15.
\end{exercice}\end{blocLivre}
\begin{blocLivre}{ex-10-3}\begin{exercice}[1]{Une médiane impaire}
Déterminer la médiane de la série : 12, 7, 9, 15, 8.
\end{exercice}\end{blocLivre}
\begin{blocLivre}{ex-10-4}\begin{exercice}[1]{Deux étendues}
Donner l’étendue de la série 12, 7, 9, 15, 8, puis celle de la série $-3$, 4, 2, $-6$.
\end{exercice}\end{blocLivre}
\begin{blocLivre}{ex-10-5}\begin{exercice}[1]{Un effectif manquant}
Une série de 25 notes est donnée par le tableau ci-dessous. Retrouver l’effectif manquant et la fréquence de la note 10.\begin{center}\renewcommand{\arraystretch}{1.35}\begin{tabular}{c|cccc}\toprule
Note & 8 & 10 & 12 & 15\\\midrule
Effectif & 4 & ? & 9 & 5\\
\bottomrule\end{tabular}\end{center}
\end{exercice}\end{blocLivre}

\rubrique{Entraînement}
\begin{blocLivre}{ex-10-6}\begin{exercice}[2]{Une moyenne pondérée}
Les valeurs 1, 4 et 6 ont pour effectifs 5, 2 et 3. Calculer la moyenne de la série.
\end{exercice}\end{blocLivre}
\begin{blocLivre}{ex-10-7}\begin{exercice}[2]{Lire un diagramme en bâtons}
Le diagramme donne le nombre de livres lus par 12 élèves pendant les vacances. Calculer la moyenne, une médiane et l’étendue.
\label{fig-ex-10-7}\figureGeo{stat-livres}
\end{exercice}\end{blocLivre}
\begin{blocLivre}{ex-10-8}\begin{exercice}[2]{Une médiane paire}
Déterminer une médiane de la série : 3, 4, 8, 10, 12, 15.
\end{exercice}\end{blocLivre}
\begin{blocLivre}{ex-10-9}\begin{exercice}[2]{Un diagramme circulaire}
Le diagramme circulaire donne la répartition des 40 élèves d’un niveau selon leur nombre de frères et sœurs. Calculer l’étendue, puis l’effectif de chaque valeur et la moyenne.
\label{fig-ex-10-9}\figureGeo{stat-fratrie}
\end{exercice}\end{blocLivre}
\begin{blocLivre}{ex-10-10}\begin{exercice}[2]{Des formules de tableur}
Les durées de trajet de dix élèves, en minutes, sont saisies en A2:A11 : 12, 15, 9, 20, 15, 11, 18, 15, 30, 10. Écrire les formules donnant la moyenne, la médiane, l’étendue et l’effectif, puis les résultats affichés.
\end{exercice}\end{blocLivre}

\rubrique{Problèmes}
\begin{blocLivre}{ex-10-11}\begin{exercice}[3]{Ajouter une valeur}
Cinq élèves ont une moyenne de 12 points à un jeu. Un sixième élève rejoint le groupe avec 18 points. Quelle est la nouvelle moyenne ?
\end{exercice}\end{blocLivre}
\begin{blocLivre}{ex-10-12}\begin{exercice}[3]{Une valeur manquante}
Les nombres 8, 11, 14 et $x$ ont une moyenne de 12. Trouver $x$, puis une médiane de la série.
\end{exercice}\end{blocLivre}
\begin{blocLivre}{ex-10-13}\begin{exercice}[3]{Une valeur extrême}
Comparer moyenne, médiane et étendue des séries A : 4, 5, 6, 7, 8 et B : 4, 5, 6, 7, 38.
\end{exercice}\end{blocLivre}
\begin{blocLivre}{ex-10-14}\begin{exercice}[3]{Comparer deux classes}
Nombre de livres empruntés dans l’année. Classe A, 10 élèves : 0, 1, 1, 2, 2, 2, 3, 5, 6, 8. Classe B : effectif 12, moyenne 3,5 livres, médiane 4 livres, étendue 6 livres.\begin{questions}
\item Quelle classe a la plus grande moyenne ?
\item Un « grand lecteur » emprunte au moins 4 livres. Quelle classe en compte le plus ?
\item Dans quelle classe les nombres de livres sont-ils les plus dispersés ?
\end{questions}
\end{exercice}\end{blocLivre}
\begin{blocLivre}{ex-10-15}\begin{exercice}[3]{Interpréter une médiane}
Dans une série de 20 durées, une médiane vaut 15 min. Peut-on conclure que dix durées valent exactement 15 min ? Donner un exemple qui justifie la réponse.
\end{exercice}\end{blocLivre}

\clearpage
